\documentstyle[%  


righttag%  


,11pt%  


,amssymb%  


,russian%               remove in Eng.  


,izvru%                 Set izven% in Eng.  


]{amsart}  


  


%  {=}


\newcommand{\const}{\operatorname{const}}  


\newcommand{\sgn}{\operatorname{sgn}}  


\newcommand{\tts}[1]{}  


\renewcommand{\baselinestretch}{0.98}


  


%\hoffset=-15mm%                  For Eng.  


  


\setcounter{page}{71}  


\god{2002}  


\nomer{4~(479)} %              Remove in Eng.  


%\nomer{Vol.\,46, No.\,4}%               Set in Eng.  


%\str{\pageref{begin}--\pageref{end}}%   Set in Eng.  


\udk{517.958:532.5}  


  


\author{\it Л.Д.\,Эскин}  


% NB:   ^^ remove in Eng.  


\title{О некоторых решениях задачи о распаде разрыва\\ в динамике  


неньютоновской жидкости}  


\thanks{Работа выполнена при финансовой поддержке Российского фонда  


фундаментальных исследований (проект \No\,00-01-00128).}  


  


\date{}  


% \pagestyle{myheadings}%         Set in Eng.  


\pagestyle{plain} %              Remove in Eng.  


\begin{document}  


% \thispagestyle{plain}%          Set in Eng.  


\maketitle  


\label{begin}  


  


  


{\bf 1.} Динамика пленочных течений нелинейно-вязкой жидкости со 


степенным реологическим законом в одномерном приближении описывается 


уравнением \cite{1}--\cite{3} 


\begin{equation} 


\frac{\partial l}{\partial t}=\frac{\partial q}{\partial x}, 


\quad t>0,\quad -\infty<x<\infty,\tag{1.1} 


\end{equation} 


где $l$ --- безразмерная толщина пленки, $q$ --- безразмерный поток 


жидкости 


\begin{equation} 


q=\sgn\frac{\partial l}{\partial x}\bigg(\frac{l^{2+n}}{n+2} 


\bigg|\frac{\partial l}{\partial x}\bigg|^n\bigg),\quad n>1,\tag{1.2} 


\end{equation} 


$t$ --- безразмерное время, постоянная $n$ определяется реологическим 


законом. 


 


Под задачей о распаде разрыва для уравнения (1.1) понимается требование 


отыскания его невозрастающего неотрицательного непрерывного на всей оси 


$x$ вместе с потоком $q$ решения $l(t,x)$, удовлетворяющего начальному 


условию 


\begin{equation} 


l(0,x)=l_0(x);\ \ l_0(x)=1\ \ (x\le0),\ \ l_0(x)=0\ \ (x>0).\tag{1.3} 


\end{equation} 


Важная роль задачи о распаде разрыва при построении решения задачи Коши 


с разрывными начальными данными хорошо известна \cite{4}. В п.\,2 


данной работы будет построено единственное автомодельное решение задачи 


(1.1)--(1.3) --- основной результат работы. В п.\,3, используя 


результаты, полученные в п.\,2, рассмотрим автомодельные разрывные 


решения уравнения (1.1), несколько модифицировав с этой целью условие 


(1.3). В дальнейшем будут использоваться обозначения 


$$ 


\alpha=\frac{1}{n+1},\ \ \beta=\frac{1}{n+2}, 


\ \ \gamma=\frac{n+1}{2n+1},\ \ \delta=\frac{1}{2n+1}, 


\ \ \lambda=\frac{1}{n-1},\ \ \mu=\frac{1}{n-2}. 


$$ 


 


{\bf 2.} В этом пункте для задачи о распаде разрыва для уравнения (1.1) 


построим невозрастающее неотрицательное непрерывное вместе с потоком 


$q=-\beta l^{n+2}(-l_x)^n$ 


решение уравнения 


\begin{equation} 


l_t=-(\beta l^{n+2}(-l_x)^n)_x,\tag{2.1} 


\end{equation} 


удовлетворяющее начальному условию (1.3). Уравнение (2.1) вместе с 


условиями (1.3) допускает группу растяжений с оператором 


\begin{equation} 


X=t\partial_t+\alpha x\partial_x.\tag{2.2} 


\end{equation} 


 


Решение задачи о распаде разрыва для уравнения (2.1) будем строить в 


виде $l(t,x)=V(\xi)$, где $\xi=x t^{-\alpha}$ --- инвариант группы 


(2.2). Из уравнения (2.1) и 


начального условия (1.3) найдем для неотрицательной убывающей функции 


$V$ нелинейное уравнение 


\begin{equation} 


\alpha\xi V_\xi=(\beta V^{n+2}(-V_\xi)^n)_\xi \tag{2.3} 


\end{equation} 


и граничные условия 


\begin{align} 


V(\infty)&=0,\tag{2.4} \\ V(-\infty)&=1.\tag{2.5} 


\end{align} 


Из условия непрерывности потока следует непрерывность функции 


$q_1=\beta V^{n+2}(-V_\xi)^n$ при 


$|\xi|<\infty$. Уравнение (2.3) рассмотрим сначала в полуплоскости 


$\xi>0$. С целью 


понижения порядка уравнения (2.3) положим 


\begin{equation} 


V=\xi^\gamma \varphi(\xi),\quad \psi=\xi \varphi'_\xi.\tag{2.6} 


\end{equation} 


Будем иметь 


\begin{equation} 


V'_\xi=\xi^{-n\delta}W(\varphi,\psi),\tag{2.7} 


\end{equation} 


$W=\psi+\gamma\varphi$. После подстановки соотношений (2.6), (2.7) в 


уравнение (2.3) 


найдем, что решения уравнения (2.3), отличные от решений $V=\const$, могут 


быть 


получены с помощью (2.6) из уравнения первого порядка 


\begin{equation} 


\frac{d\psi}{d\varphi}= 


-\frac{P_1(\varphi,\psi)}{Q_1(\varphi,\psi)}, \tag{2.8} 


\end{equation} 


где 


$$ 


P_1=-\frac{\alpha}{\beta}+\varphi^{n+1}(-W)^{n-2} 


[\gamma n\varphi\psi+\delta(3n+2)\varphi W+\tfrac{1}{\beta} 


\psi W],\quad Q_1=n\varphi^{n+2}\psi(-W)^{n-2}. 


$$ 


Каждое решение $\psi=\psi(\varphi)$ уравнения (2.8) порождает 


однопараметрическое 


семейство решений уравнения с разделяющимися переменными 


$\psi=\xi\varphi'_\xi$, а затем с 


помощью первого из соотношений (2.6) и однопараметрическое семейство 


решений уравнения (2.3). 


 


Прежде всего построим в полуплоскости $\xi>0$ невозрастающее решение 


$V_1$ 


уравнения (2.3), удовлетворяющее граничному условию (2.4) и условию 


непрерывности потока $q$. В силу (2.7) и условия невозрастания решения 


$V_1$ уравнение (2.8) следует рассматривать в секторе I, являющемся 


пересечением полуплоскости $\varphi>0$ с полуплоскостью $W<0$. Нетрудно 


заметить, 


что сектору I целиком принадлежит непрерывная ветвь 


$\psi=\psi_1(\varphi)$ изоклины нуля 


(кривая $P_1(\varphi,\psi)=0$) уравнения (2.8). 


 


\begin{center}


\unitlength=1mm


\begin{picture}(170,87) %% width, heigth, added 5 mm for signature


\put(42,87){\special {em:graph es1.bmp}} %% left X, upper Y, -, /2, +toX


\put(78,0){\mbox Рис.\,1}


\end{picture}


\end{center} 





 


\noindent{}Эта ветвь при $\varphi\to0$ имеет вертикальную асимптоту 


$\varphi=0$, а при $\varphi\to\infty$ 


--- наклонную асимптоту $\psi=k_1\varphi$, где угловой коэффициент 


$$ 


k_1=-\frac{\beta\delta}{2}(2n^2+7n+4+(4n^4+16n^3+21n^2+8n)^{1/2}), 


$$ 


причем $\psi_1\sim-(\alpha\varphi^{-n-1})^{1/n}$ при $\varphi\to0$. 


Используя этот факт, можно показать, что в секторе~I


имеется два семейства интегральных кривых уравнения (2.8). Кривые 


семейства $L_1$ начинаются (рис.\,1) в точках прямой $A$: $W=0$. При 


$\varphi\to\infty$ они имеют 


прямую $A$ своей наклонной асимптотой, причем справедлива асимптотика 


\begin{equation} 


W(\varphi,\psi)\sim B_1\varphi^{-n\alpha}.\tag{2.9} 


\end{equation} 


В соотношении (2.9) коэффициент $B_1=B_1(\varphi_0)<0$ ($\varphi_0>0$) 


однозначно определяется точкой 


$(\varphi_0,\psi_0)$, в которой начинается интегральная кривая семейства 


$L_1$. Эти кривые 


либо не имеют экстремальных точек, либо имеют по две экстремальные 


точки (точку максимума и точку минимума), расположенные на изоклине 


нуля $\psi_1$. Интегральные кривые семейства $L_2$ имеют вертикальную 


асимптоту 


$\varphi=0$ и наклонную асимптоту --- прямую $A$, причем для них 


справедливы 


асимптотики 


\begin{equation} 


\psi\sim A_1\varphi^{-1/(n\beta)}\ \ (\varphi\to0),\quad 


W(\varphi,\psi)\sim c_1(A_1)\varphi^{-n\alpha}\ \ (\varphi\to\infty). 


\tag{2.10} 


\end{equation} 


Здесь коэффициенты $A_1<0$ и $c_1=c_1(A_1)<0$ определяются интегральной 


кривой семейства 


$L_2$, всегда $B_1(\varphi_0)>c_1(A_1)$. Кривые семейства $L_2$ 


пересекают изоклину нуля $\psi_1$ в точке 


максимума. Семейства $L_1$ и $L_2$ разделяет также лежащая в секторе I 


кривая 


$L$ --- интегральная кривая с асимптотикой 


\begin{alignat}{2} 


\psi&\sim -\bigg(\frac{\beta}{\alpha}\varphi^{n+1} 


\bigg)^{-1/n} &\quad &(\varphi\to0), \tag{2.11} \\ 


W(\varphi,\psi)&\sim B\varphi^{-n\alpha} &\quad 


&(\varphi\to\infty).\tag{2.12} 


\end{alignat} 


Коэффициент $B$ в соотношении (2.12) может быть определен лишь численно. 


Отметим, что всегда (т.\,е. для всех $\varphi_0>0$ и $A_1<0$) справедливы 


неравенства 


$B_1(\varphi_0)>B>c_1(A_1)$. 


 


Однопараметрические семейства $L_1$ и $L_2$ порождают в силу соотношений 


(2.6), (2.7) двупараметрические семейства неотрицательных убывающих при 


$\xi>0$ решений уравнения (2.3), кривая $L$, в свою очередь, порождает 


однопараметрическое семейство решений (параметр семейства --- точка 


фронта $\xi_0$ решения $V(\xi)$). С помощью асимптотических соотношений 


(2.9), 


(2.10) нетрудно убедиться, что решения, порождаемые в полуплоскости 


$\xi>0$ 


интегральными кривыми семейств $L_1$ и $L_2$, не могут удовлетворять 


граничному условию (2.4) и условию непрерывности потока $q$ при $\xi>0$. 


 


Наоборот, из соотношения (2.11) следует 


$$ 


\varphi\sim c_n\bigg(\frac{\xi_0-\xi}{\xi_0} \bigg)^{n\delta} 


\ \ (\xi\to\xi_0-0,\ \; \xi_0>0),\ \ c_n= 


(n\delta)^{-n\delta}(\alpha/\beta)^\delta, 


$$ 


откуда для неотрицательных убывающих при $\xi>0$ решений $V_1$ уравнения 


(2.3), 


порождаемых кривой $L$, получаем слева от точки фронта $\xi_0$, которая 


пока 


остается произвольной, асимптотические представления 


\begin{equation} 


V_1\sim c_n\xi^\delta_0(\xi_0-\xi)^{n\delta} 


\ \ (\xi\to\xi_0-0).\tag{2.13} 


\end{equation} 


С помощью соотношений (2.7) и (2.11) находим асимптотику производной 


\begin{equation} 


V'_1\sim -n\delta c_n\xi^\delta_0(\xi_0-\xi)^{-\gamma} 


\ \ (\xi\to\xi_0-0).\tag{2.14} 


\end{equation} 


Из (2.13), (2.14) без труда следует, что в точке фронта $\xi_0$ решения 


$V_1$ 


поток $q=0$, поэтому, продолжив построенное на отрезке $0\le\xi\le\xi_0$ 


решение $V_1$ с 


помощью соотношения $V_1(\xi)=0$ при $\xi>\xi_0$, получим определенное на 


полуоси $\xi\ge0$ 


неотрицательное невозрастающее решение уравнения (2.3), удовлетворяющее 


граничному условию (2.4) и условию непрерывности потока при $\xi>0$. 


 


Воспользуемся теперь асимптотикой (2.12). Найдем, что $\varphi\sim 


V_{10}\xi^{-\gamma}$ ($\xi\to+0$), откуда в 


силу (2.6), (2.7) получим 


\begin{equation} 


V_1(0)=V_{10},\quad V'_1(0)=B V^{-n\alpha}_{10}.\tag{2.15} 


\end{equation} 


Постоянная $V_{10}$ в (2.15) взаимно однозначно определяется точкой 


фронта $\xi_0$ 


решения $V_1(\xi)$. 


 


Чтобы продолжить построенное решение $V_1$ на полуплоскость $\xi<0$, 


положим 


аналогично (2.6) 


\begin{equation} 


V(\xi)=(-\xi)^\gamma \varphi(\xi),\quad 


\psi=\xi\varphi'_\xi,\ \ \xi<0,\tag{2.16} 


\end{equation} 


откуда найдем 


\begin{equation} 


V'_\xi=-(-\xi)^{n\delta}W(\varphi,\psi).\tag{2.17} 


\end{equation} 


После подстановки соотношений (2.16), (2.17) в уравнение (2.3) найдем 


\begin{equation} 


\frac{d\psi}{d\varphi}=-\frac{P_2(\varphi,\psi)} 


{Q_2(\varphi,\psi)},\tag{2.18} 


\end{equation} 


где 


$$ 


P_2=\frac{\alpha}{\beta}+\varphi^{n+1}W^{n-2} 


[\gamma n\varphi\psi+\delta(3n+2)\varphi W+\tfrac{1}{\beta} 


\psi W],\quad Q_2=n\varphi^{n+2}\psi W^{n-2}. 


$$ 





\begin{center}


\unitlength=1mm


\begin{picture}(170,127) %% width, heigth, added 5 mm for signature


\put(35,127){\special {em:graph es2.bmp}} %% left X, upper Y, -, /2, +toX


\put(78,0){\mbox Рис.\,2}


\end{picture}


\end{center} 





\noindent{}Неотрицательные убывающие в полуплоскости $\xi<0$ решения $V(\xi)$ 


уравнения (2.3) 


в силу соотношений (2.16), (2.17) порождаются интегральными кривыми 


уравнения (2.18), принадлежащими сектору II плоскости $\varphi$, $\psi$, 


который 


является пересечением полуплоскостей $\varphi>0$ и $W>0$ (рис.\,2). 


 


Придется по отдельности рассматривать три случая: 1)~$n>2$; 


2)~$n=2$; 3)~$1<n<2$. Во всех трех случаях изоклина нуля уравнения (2.19) 


имеет 


в секторе II ветвь $\psi_2$ с наклонной асимптотой $\psi=k_2\varphi$, где 


$$ 


k_2=-\frac{\beta\delta}{2}(2n^2+7n+4-(4n^4+16n^3+21n^2+8n)^{1/2}). 


$$ 


В случае 1) изоклина нуля имеет в секторе II еще и вторую ветвь $\psi_3$ с 


наклонной асимптотой $W=0$, причем справедлива асимптотика 


$$ 


\psi_3\sim -\gamma\varphi+c\varphi^{-(n+3)\mu},\quad 


c=\bigg(\frac{\alpha^3}{n\beta\delta^2} \bigg)^\mu 


\ \ (\varphi\to\infty). 


$$ 


Ветви $\psi_2$ и $\psi_3$ не могут пересекаться с границей сектора II и с 


осью $\varphi$ и 


соединяются между собой в некоторой внутренней точке сектора II. В 


случае~2) ветвь $\psi_3$ исчезает, а ветвь $\psi_2$ пересекает границу 


сектора II в 


точке $((50/27)^{1/5},-3(50/27)^{1/5}/5)$, принадлежащей прямой $W=0$. В 


случае~3), переписав уравнение 


(2.18) в эквивалентной форме $\psi'_\varphi=-P_3/Q_3$ ($P_3=P_2W^{2-n}$, 


$Q_3=n\varphi^{n+2}\psi$), снова получаем, что ветвь $\psi_3$ 


исчезает, а ветвь $\psi_2$ пересекает прямую $W=0$ в начале координат. 


Используя эти результаты, можно показать, исследуя поведение интегральных


кривых, что в каждом из трех случаев в секторе II существуют


два семейства интегральных 


кривых уравнения (2.18) (рис.\,2). Кривые обоих семейств $M_1$ и $M_2$ 


начинаются в точках $(\varphi_0,\psi_0)$ прямой $W=0$, но кривые 


семейства $M_1$ имеют 


вертикальную асимптоту $\varphi=0$, а кривые семейства $M_2$ --- наклонную 


асимптоту $W=0$. Кривые семейства $M_1$ в случае~1) либо не пересекают ни 


одной ветви изоклины нуля, либо пересекают обе ветви $\psi_2$ и $\psi_3$ 


(ветвь $\psi_2$ 


интегральная кривая пересекает в точке минимума, ветвь $\psi_3$ --- в 


точке 


максимума). Кривые семейства $M_2$ пересекают лишь ветвь $\psi_3$, а 


затем, 


убывая при $\varphi\to\infty$, приближаются к асимптоте $W=0$. В 


случаях~2) и~3) кривые 


семейства $M_2$ не пересекают ветвь $\psi_2$ изоклины нуля. Кривые 


семейства $M_1$ 


в случае~2) могут как пересекать, так и не пересекать ветвь $\psi_2$, в 


случае~3) все кривые семейства $M_1$ пересекают ветвь $\psi_2$. Семейства 


$M_1$ и 


$M_2$ разделены интегральной кривой $M$, выходящей из некоторой точки 


$(\wt\varphi_0,\wt\psi_0)$ 


прямой $W=0$ и уходящей в секторе II в особую точку на $\infty$ в 


направлении с 


угловым коэффициентом 


\begin{equation} 


k_3=-(3n+2)\alpha\delta/2,\quad k_2>k_3>-\gamma.\tag{2.19} 


\end{equation} 


Для кривых семейства $M_1$\; $\varphi_0<\wt\varphi_0$, для кривых 


семейства $M_2$\; $\varphi_0>\wt\varphi_0$. Во всех 


трех случаях справедлива асимптотика для кривых семейства $M_1$ 


\begin{equation} 


\psi\sim A_2\varphi^{-1/n\beta}\ \ (\varphi\to0),\tag{2.20} 


\end{equation} 


аналогичная первой из формул (2.10). Для кривых семейства $M_2$ также 


справедлива асимптотика 


\begin{equation} 


W(\varphi,\psi)\sim B_2\varphi^{-n\alpha}\ \ (\varphi\to0),\tag{2.21} 


\end{equation} 


аналогичная второму соотношению в (2.10). Коэффициенты 


$A_2=A_2(\varphi_0)>0$ и $B_2=B_2(\varphi_0)>0$ в 


формулах (2.20), (2.21) теперь однозначно определяются абсциссой 


$\varphi_0$ 


точки прямой $W{=}0$, из которой выходит интегральная кривая, причем с 


возрастанием $\varphi_0$ коэффициент $A_2(\varphi_0)$ возрастает, 


коэффициент $B_2(\varphi_0)$ непрерывно 


убывает от $\infty$ до 0 при возрастании $\varphi_0$ от $\wt\varphi_0$ до 


$\infty$. 


 


Исследуем поведение интегральных кривых семейств $M_1$ и $M_2$ вблизи 


точки 


выхода $(\varphi_0,-\gamma\varphi_0)$, принадлежащей прямой $W=0$, т.\,е. 


при $\varphi\to\varphi_0$, $\psi\to\psi_0=-\gamma\varphi_0$. Получим 


\begin{gather} 


\psi-\psi_0=d_n(\varphi-\varphi_0)^\lambda-\gamma 


(\varphi-\varphi_0)+O((\varphi-\varphi_0)^{n\lambda}) 


\ \ (\varphi\to\varphi_0-0), \tag{2.22} \\ 


d_n=\bigg(\frac{\alpha^2}{n\beta\delta\lambda\varphi^{n+3}_0} 


\bigg)^\lambda.\notag 


\end{gather} 


С помощью соотношений (2.16), (2.17), (2.19) и асимптотических 


представлений (2.20)--(2.22) теперь нетрудно исследовать поведение 


двупараметрических семейств $N_1$ и $N_2$ неотрицательных монотонно 


убывающих 


решений уравнения (2.3), порождаемых соответственно 


однопараметрическими семействами $M_1$ и $M_2$, а также поведение 


однопараметрического семейства $N$ неотрицательных монотонно убывающих 


решений, порождаемых интегральной кривой $M$. Оказывается, что все 


решения семейства $N$ обращаются в нуль при $\xi=0$, а все решения 


семейства 


$N_1$ пересекают отрицательную полуось $\xi<0$ (каждое в соответствующей 


точке). Следовательно, решения, принадлежащие этим семействам, 


невозможно гладко срастить в точке $\xi=0$ с решением $V_1$. Поэтому 


интересны 


лишь решения семейства $N_2$. Эти решения определены при $0>\xi>\xi_1$ 


($\xi_1$ и $\varphi_0$ --- 


два параметра, определяющие решение, параметр $\xi_1$ выбирается так, 


чтобы 


$V'_2(\xi_1)=0$, возможность такого выбора обусловлена справедливостью 


асимптотики 


(2.22)). В результате достаточно громоздких выкладок, которые опустим, 


с помощью асимптотики (2.22) получим для решений семейства $N_2$ при 


$\xi\to\xi_1+0$ 


\begin{align} 


V_2&=(-\xi_1)^\gamma \varphi_0\bigg[1- 


\frac{\delta d_n}{n\alpha\lambda(-\psi_0)^{\lambda/\mu}} 


\bigg(\frac{\xi_1-\xi}{\xi_1}\bigg)^{n\lambda}+ 


O((\xi-\xi_1)^w)\bigg],\quad w=\min\{2,\lambda/\alpha \},\notag \\ 


% 


V_2&=(-\xi_1)^{3/5}\varphi_0\bigg[1+\frac{3}{5}\bigg( 


\frac{d_2}{2}-\frac{8}{5}\bigg)\bigg(\frac{\xi_1-\xi}{\xi_1} 


\bigg)^2+O((\xi-\xi_1)^3)\bigg],\quad n=2, \tag{2.23} \\ 


% 


V_2&=(-\xi_1)^\gamma \varphi_0[1+o((\xi-\xi_1)^2)],\quad 


1<n<2.\notag 


\end{align} 


Из (2.17) с помощью (2.22) при $\xi\to\xi_1+0$ для всех $n>1$ найдем 


\begin{equation} 


V'_2=-(-\xi_1)^{-n\delta}d_n(\gamma\varphi_0)^\lambda 


\bigg(\frac{\xi_1-\xi}{\xi_1} \bigg)^\lambda+ 


O((\xi-\xi_1)^{n\lambda}).\tag{2.24} 


\end{equation} 


 


С другой стороны, с помощью асимптотики (2.21) нетрудно убедиться, что 


для решений семейства $N_2$ при $\xi\to-0$ существует конечный $\lim 


V_2(\xi)=V_{20}$, зависящий от $\xi_1$ 


и $\varphi_0$, причем $V'_2(0)=-B_2V^{-n\alpha}_{20}$, 


$B_2=B_2(\varphi_0)$. 


 


Теперь, положив $\xi_1=-\varphi^{1/\gamma}_0$, в силу соотношений 


(2.23) и (2.24) получим решение 


семейства $N_2$, удовлетворяющее условиям $V_2(\xi_1)=1$ и 


$V'_2(\xi_1)=0$. Продолжим его для 


значений $\xi<\xi_1$, положив $V_2(\xi)=1$. В результате в 


полуплоскости $\xi<0$ получим 


решение $V_2(\xi)$, удовлетворяющее граничному условию (2.5) и условию 


непрерывности потока при $\xi<0$. Это решение зависит от параметра 


$\varphi_0$, 


который выберем так, чтобы иметь 


\begin{equation} 


B_2(\varphi_0)=-B.\tag{2.25} 


\end{equation} 


Напомним, что $B_2(\varphi_0)$ непрерывно убывает от $\infty$ до 0 с 


возрастанием $\varphi_0$ от $\wt\varphi_0$ 


до $\infty$, а $B<0$, откуда следует существование единственного 


решения $\varphi_0$ 


уравнения (2.25). 


 


Наконец, выбор точки фронта $\xi_0>0$ подчиним требованию 


\begin{equation} 


V_{20}=V_{10}.\tag{2.26} 


\end{equation} 


Существование единственного решения уравнения (2.26) следует из легко 


проверяемого группового свойства решений уравнения (2.3): вместе с 


решением $V(\xi)$ решением является и функция 


$\zeta^{-\gamma}V(\zeta\xi)$, $\zeta>0$ произвольно. В самом 


деле, рассмотрим порожденное кривой $L$ решение $\wt V(\xi)$ уравнения 


(2.3) с 


точкой фронта $\xi_0=1$ (оно определено однозначно) и положим 


\begin{equation} 


V_1(\xi)=\zeta^{-\gamma} \wt V(\zeta\xi).\tag{2.27} 


\end{equation} 


С помощью уравнения (2.26) без труда найдем, что при $\zeta=(V_{20}/\wt 


V(0))^{-1/\gamma}$ определенное 


равенством (2.27) решение $V_1$ уравнения (2.3) с точкой фронта 


$\xi_0=1/\zeta$ и 


нулевым потоком $q_1$ на фронте удовлетворяет как соотношению (2.26), так 


и равенству 


\begin{equation} 


V'_1(0)=V'_2(0),\tag{2.28} 


\end{equation} 


т.\,е. является искомым. Положив 


$$ 


V(\xi)= 


\begin{cases} 


V_1(\xi), & \xi\ge0; \\ 


V_2(\xi), & \xi<0, 


\end{cases} 


$$ 


получим определенное на всей оси решение уравнения (2.3), 


удовлетворяющее как граничным условиям (2.4), (2.5), так и условию 


непрерывности потока $q_1$, что следует из (2.15), (2.25), (2.26), (2.28). 


 


Предложенный выше алгоритм решения задачи о распаде разрыва и 


полученные результаты удобно сформулировать в следующем виде. 


\medskip 


 


{\bf Теорема.} {\it Существует единственное невозрастающее и 


непрерывное вместе с потоком $q$ $(1.2)$, неотрицательное автомодельное 


решение $l(t,x)=V(\xi)$, $\xi=x t^{-\alpha}$ задачи о распаде разрыва 


$(1.1)$, $(1.3)$.} 


\medskip 


 


Это решение $V(\xi)=1$ при $\xi\le\xi_1=-\varphi^{-1/\gamma}_0$, где 


$\varphi_0>0$ --- единственное решение уравнения 


(2.25), в котором $B<0$ --- коэффициент в асимптотике (2.12), а 


$B_2(\varphi_0)>0$ --- 


коэффициент в асимптотике (2.21). При $\xi_1<\xi<0$ решение $V(\xi)$ 


совпадает с 


решением $V_2(\xi)$ обыкновенного дифференциального уравнения (2.3), для 


которого при $\xi\to\xi_1+0$ справедливы асимптотические представления 


(2.23), 


(2.24), так что $V_2(\xi_1)=1$, $V'_2(\xi_1)=0$. При 


$0\le\xi\le\zeta^{-1}$ решение $V(\xi)=\zeta^{-\gamma}\wt V(\zeta\xi)$, 


здесь $\wt V(\xi)$ --- решение 


уравнения (2.3) с точкой фронта $\xi_0=1$ и асимптотикой (2.13), (2.14) 


при $\xi\to1-0$ 


(оно определяется единственным образом), а $\zeta=(V_2(0)/\wt 


V(0))^{-1/\gamma}$. Наконец, $V(\xi)=0$ при $\xi>\zeta^{-1}$. 


\medskip 


 


{\bf 3.} Применим полученные в п.\,2 результаты для построения 


автомодельных разрывных неотрицательных невозрастающих решений 


модифицированной задачи (1.1), (1.3). С этой целью прежде всего выведем 


необходимое условие, которое должно выполняться для разрывного решения 


$l$ на линии разрыва $x=x(t)$. Обозначим 


\begin{alignat*}{3} 


l^+&=\lim_{x\to x(t)+0}l(t,x), 


&\ \ l^-&=\lim_{x\to x(t)-0}l(t,x), 


&\ \ [l]&=l^+-l^-, \\ 


q^+&=\lim_{x\to x(t)+0}q, &\ \ q^-&=\lim_{x\to x(t)-0}q, 


&\ \ [q]&=q^+-q^-, 


\end{alignat*} 


$x_0(t)$ --- точка фронта решения $l$ ($l(t,x_0(t))=0$, 


$q|_{x=x_0(t)}=0$), $x_1(t)$ --- точка сращивания 


криволинейного участка решения $l(t,x)$ с прямой $l=1$. Из закона 


сохранения 


массы находим равенство 


\begin{equation} 


\int\limits^{x(t)-0}_{x_1(t)}l(t,x)d x+ 


\int\limits^{x_0(t)}_{x(t)+0}l(t,x)d x=-x_1(t).\tag{3.1} 


\end{equation} 


Дифференцируя тождество (3.1) по $t$, получаем 


\begin{equation} 


x'(l^--l^+)+\int\limits^{x(t)-0}_{x_1(t)}l_t d x+ 


\int\limits^{x_0(t)}_{x(t)+0}l_t d x=0.\tag{3.2} 


\end{equation} 


С учетом уравнения (1.1) из (3.2) следует условие 


\begin{equation} 


x'(t)[l]=-[q],\tag{3.3} 


\end{equation} 


которое и должно выполняться на линии разрыва $x=x(t)$. 


 


Рассмотрим в плоскости $x$, $t$ область $D$, ограниченную снизу 


полупрямой 


$R_1=\{(x,t):x<0,\ t=0\}$, а справа --- кривой 


\begin{equation} 


R_2=\{(x,t):x=\xi_0t^\alpha,\ \xi_0>0 \text{ задано }\}.\tag{3.4} 


\end{equation} 


В области $D$ построим разрывное решение $l$ уравнения (1.1), 


удовлетворяющее следующим условиям на границе области 


\begin{gather} 


l(t,x)=1,\quad (x,t)\in R_1,\tag{3.5} \\ 


l(t,x)=0,\ \ q=0,\ \ (x,t)\in R_2.\tag{3.6} 


\end{gather} 


Кроме того, на линии разрыва 


\begin{equation} 


R_3=\{(x,t):x=\xi_2t^\alpha,\ \xi_2<0 \text{ задано }\} \tag{3.7} 


\end{equation} 


в соответствии с (3.3) потребуем выполнения соотношения 


\begin{equation} 


\alpha\xi_2t^{\alpha-1}[l]=[\wt q],\quad 


\wt q=-q=\beta l^{n+2}(-l_x)^n.\tag{3.8} 


\end{equation} 


Будем рассматривать обобщенное решение задачи (1.1), (3.5)--(3.8), 


допуская, как обычно \cite{5}, что искомое решение может не только 


иметь разрыв первого рода на кривой $R_3$, но может не иметь всех входящих 


в уравнение (1.1) производных и на некоторых других кривых в области~$D$. 


\medskip 


 


{\bf Замечание.} Отметим наличие разрыва в граничных данных 


сформулированной задачи. Действительно, из (3.5) находим 


$\lim l=1$ ($x\to-0$, $(x,t)\in R_1$), а из (3.6) $\lim l=0$ ($x\to+0$, 


$(x,t)\in R_2$). 


\medskip 


 


Снова будем строить автомодельное неотрицательное невозрастающее (в 


точках непрерывности) решение задачи (1.1), (3.5)--(3.8) (этой цели и 


отвечает выбор кривых $R_2$ и $R_3$ в виде (3.4) и (3.7)), положив 


$l(t,x)=V(\xi)$, $\xi=x t^{-\alpha}$. Из 


(1.1), (3.5)--(3.8) следует, что функция $V(\xi)$ должна удовлетворять 


уравнению (2.3) при $\xi<\xi_0$ всюду, кроме точки разрыва $\xi_2$, в 


которой в силу 


(3.7), (3.8) должно выполняться соотношение 


\begin{equation} 


\alpha\xi_2[V]=[q_1],\quad q_1=\beta V^{n+2}(-V_\xi)^n\tag{3.9} 


\end{equation} 


(в некоторых точках полуоси $\xi<\xi_0$ функция $V$ может не иметь всех 


производных, входящих в уравнение (2.3)). Граничные условия (3.5), 


(3.6) дают 


\begin{alignat}{2} 


V&=1 &\quad (\xi&=-\infty),\tag{3.10}\\ 


V&=q_1=0 &\quad (\xi&=\xi_0).\tag{3.11} 


\end{alignat} 


Наконец, функция $V$ должна быть неотрицательной и невозрастающей. 


 


Перейдем к построению искомого решения $V(\xi)$ уравнения (2.3), 


удовлетворяющего условиям (3.9)--(3.11). При $0\le\xi\le\xi_0$ полагаем 


$V(\xi)=V_1(\xi)$, где $V_1$ --- 


построенное в п.\,2 решение уравнения (2.3), удовлетворяющее условию 


(3.11). Оно определено однозначно, поскольку точка фронта $\xi_0>0$ 


задана, 


неотрицательнo и не возрастает. Таким образом, определены и значения 


$V_1(0)>0$ и $V'_1(0)<0$. Теперь при $\xi<0$ можно построить


единственное неотрицательное невозрастающее решение 


$V_2(\xi)$, удовлетворяющее условиям $V_2(0)=V_1(0)$, $V'_2(0)=V'_1(0)$. 


Это решение принадлежит семейству 


$N_2$ и метод его построения указан в п.\,2. 


 


Вместе с решением $V_2$ оказываются определенными значения $V^+_2=\lim 


V_2$, $q^+_1=\lim q_1$, $\xi\to\xi_2+0$, 


следовательно, определено и значение $m^+=\alpha\xi_2V^+_2-q^+_1$ (т.\,к. 


$\xi_2<0$, $V^+_2>0$, $q^+_1>0$). Полагаем 


$V(\xi)=V_2(\xi)$, $\xi_2<\xi<0$. 


 


При $\xi<\xi_2$ построим решение $V_3(\xi,\xi_1,\varphi_0)$ уравнения 


(2.3), принадлежащее семейству $N_1$, 


если $0<\varphi_0<\wt\varphi_0$, семейству $N$, если 


$\varphi_0=\wt\varphi_0$, семейству $N_2$ в случае 


$\varphi_0>\wt\varphi_0$ (см. п.\,2). 


Будем иметь $V_3=(-\xi_1)^\gamma \varphi_0$, $V'_3=0$ при $\xi=\xi_1$. 


Положив $\xi_1=-\varphi^{-1/\gamma}_0<0$, получим $V_3=1$, $V'_3=0$ при 


$\xi=\xi_1$. Теперь 


оказываются определенными значения 


$V^-_3(\varphi_0)=\lim\limits_{\xi\to\xi_2-0} 


V_3(\xi,\xi_1,\varphi_0)>0$, $q^-_1=\lim\limits_{\xi\to\xi_2-0} \beta 


V^{n+2}_3(-V'_3)^n>0$, а вместе с ними и значение 


$m^-(\varphi_0)=\alpha\xi_2V^-_3-q^-_1<0$. Соотношение (3.9) равносильно 


трансцендентному уравнению 


\begin{equation} 


m^-(\varphi_0)=m^+ \tag{3.12} 


\end{equation} 


относительно $\varphi_0>0$. После решения уравнения (3.12) определится 


функция $V_3(\xi,\xi_1,\varphi_0)$ 


($\xi_1=-\varphi^{1/\gamma}_0$). Положив $V=V_3$ при $\xi_1<\xi<\xi_2$, 


$V=1$ при $\xi<\xi_1$, получим искомое решение $V(\xi)$. 


 


Для доказательства разрешимости уравнения (3.12) рассмотрим поведение 


решения $V_3$ и его производной $V'_3$ при изменении $\varphi_0$ от 0 до 


$\infty$. Если $\varphi_0<\wt\varphi_0$, 


то кривая $V_3$ принадлежит семейству $N_1$ и при $\varphi_0\to0$ будем 


иметь $\varphi(\xi)\to0$ 


равномерно по $\xi$, а в силу (2.20)\; $W(\varphi,\psi)\sim 


A_2(\varphi_0)\varphi^{-1/n\beta}$. Отсюда и из (2.16), (2.17) 


получим при $\varphi_0\to0$ 


$$ 


V^-_3\to0,\ \ q^-_1\sim \beta(-\xi_2)^{n+1+\delta} 


A^n_2(\varphi_0)\to0, 


$$ 


т.\,к. $A_2(\varphi_0)\to0$ при $\varphi_0\to0$. Следовательно, 


$m^-(\varphi_0)\to0$ при $\varphi_0\to0$. При $\varphi_0>\wt\varphi_0$ 


кривая $V_3$ 


принадлежит семейству $N_2$ и при $\varphi_0\to\infty$ будем иметь 


$\xi_1\to-0$. Отсюда найдем $V^-_3=1$ 


(т.\,к. $\xi_2<\xi_1$, если $\varphi_0$ достаточно велико), $q^-_1=0$ 


($V'_3(\xi_2-0)=0$, $\varphi_0\gg1$). Итак, имеем $m^-=\alpha\xi_2<0$, 


если $\varphi_0\gg1$. Таким образом, значение $m^-(\varphi_0)$ при 


изменении $\varphi_0$ от 0 до $\infty$ 


принимает, во всяком случае, все отрицательные значения межу нулем и 


$\alpha\xi_2$ 


(нетрудно заметить, что $m^-(\varphi_0)$ --- непрерывная функция 


$\varphi_0$). Следовательно, 


достаточным условием разрешимости уравнения (3.12), а вместе с ним и 


достаточным условием существования автомодельного неотрицательного и 


невозрастающего решения $l=V(\xi)$ смешанной задачи (1.1), (3.5)--(3.8) в 


неограниченной области $D$ будет выполнение неравенства 


$m^+>\alpha\xi_2$, что можно 


рассматривать как определенное ограничение на выбор значения параметра 


$\xi_0$ (или $\xi_2$). Отметим, что построенное решение $l(t,x)$ помимо 


разрыва на 


кривой $R_3$ допускает и разрыв второй производной по $x$ на кривой 


$x=\xi_1t^\alpha$, $t>0$. 


 


Искренне благодарю А.Н.\,Саламатина и В.А.\,Чугунова за полезные 


обсуждения. 


 


 


\begin{thebibliography}{9} 


 


\bibitem{1} 


Ritz C. {\em Un modele termo-mecanique d'evolution pour le bassin


glaciaire antarctique Vostok-Glacier Byrd$:$ sensibilite aux valeurs des


parametres mal connus} // These Doctoral d'Etat. Centre national de la


recherhe scientifique. Laboratorie de Glaciologie et Geophisique de


l'Environment, 1992. -- 377 p.


\bibitem{2} 


Саламатин А.Н., Чугунов В.А., Мазо А.Б. {\em Численные и инвариантные 


решения задачи о динамике субизотермического ледника в одномерном 


приближении} // Задачи механики природных процессов. -- М.: Изд-во МГУ, 


1983. -- С.\,82--96. 


\bibitem{3} 


Тонконог С.Л., Чугунов В.А., Эскин Л.Д. {\em Методы группового анализа 


в задаче о растекании с проскальзыванием нелинейно вязкой жидкости} // 


ПММ. -- 1994. -- Т.\,58. -- \No\,4. -- С.\,63--69. 


\bibitem{4} 


Гельфанд И.М. {\em Некоторые задачи теории квазилинейных уравнений} // 


УМН. -- 1959. -- Т.\,14. -- \No\,2. -- С.\,87--158. 


\bibitem{5} 


Самарский А.А., Галактионов В.А., Курдюмов С.П., Михайлов А.П. {\em 


Режимы с обострением в задачах для квазилинейных параболических 


уравнений}. -- М.: Наука, 1987. -- 480 с. 


 


\end{thebibliography} 


 


\vskip 2cm 


 


\postup{\it Казанский государственный}{\it Поступили} 


\postup{\it университет}{\it первый вариант $22.09.1999$} 


\postup{}{\it окончательный вариант $09.10.2000$} 


 


\label{end} 


\end{document} 


